\section{Summing the remainder at arrival}\label{app:arrival}
A single geometric square can contribute on the scale \(1/n\) relevant to \eqref{eq:signed-mean}. After we
subtract \(\Gamma\), each surviving support must satisfy the geometric constraints of two geometric squares.
This additional cost lets us sum the errors with a bound that decreases with \(b\). Consider first the
short historical arrival event \(H_j\) under \(\nu\). For a set of board squares \(Y\), we activate only the
diagonal constraints supported on \(Y\), keeping the entire rook mask. Let \(e(Y)\) be the logarithm
of the avoidance probability minus the change in the predictor, and let \(e_C(Y)\) be the same
expression using only the pairs of \(C\). The coefficient of a support \(U\) is \(\sum_{Y\subseteq
U}(-1)^{|U|-|Y|}e(Y)\). An isolated board square cancels when we pair subsets that differ only in that
square. If all pairs belong to one geometric square, the coefficient cancels against that of \(e_C\). If two
distinct geometric squares are involved, no \(e_C\) contributes, since the board squares would have to lie
at the four vertices of \(C\), whose only diagonal pairs are the two of \(C\). The remainder therefore
sums the original coefficients of supports that have no isolated board squares and contain at least two
distinct geometric squares. The subtracted part has degree at most four.

Let \(v\) be the number of board squares, \(h\) the total number of distinct row and column labels, \(D\)
the number of free coordinates, and \(p\) the support exponent. Prescribing the \(h\) lines gives the
factor \((j/n)^h\), together with \(\delta^{-2v}\exp(C_Lv^2/j)\) in the prescription bound. For each
type, the number of positions is of order \(n^D\). Write \(d_*=h-D\), \(q=p-D-1\). Two distinct
geometric squares impose two nonproportional linear equations on the row and column coordinates. Hence
\(d_*\ge2\), and the support bound gives \(q\ge0\). The scaled factor is
\[
 n\,n^D j^{-p}(j/n)^h=(j/n)^{d_*-1}j^{-q}
 \le (b/\delta_*)\delta_*^{d_*-1}j^{-q}
 \qquad(0<b\le\delta_*).
\]
We factor out \(b/\delta_*\) before summing. For \(v\le4\) there is an absolute bound on the number of
types. For noncollinear types the weighted count uses \(d_*-1+q\ge(v-4)/3\). Thus for \(5\le v\le
N=\lceil4\log n/\log\rho\rceil\), with \(\rho\ge2\) fixed, the bound after summing over types of that
size is
\[
 (b/\delta_*) C_L(C_1\delta^{-2})^v v^{d_{\rm pol}}\delta_*^{(v-4)/3}.
\]
Choose \(\delta_*\) so that \(C_1\delta^{-2}\delta_*^{1/3}<1\). For a collinear type with \(v\ge3\) we
have \(p=v+2,h=2v\), with at most \(4n\,n^v/v!\) supports. This gives \(C(C_2\delta^{-2})^v
v^2(j/n)^{v-2}\), and we decrease \(\delta_*\) once to sum this bound as well. Substituting atoms up
to degree \(N\) incurs an error of \(O(j^{-8})\) in value, and the global norm bounds the tail by
\(O_\rho(n^{-2})\). Both errors vanish after scaling for each fixed \(b\). Thus \(\limsup_n n\mathbb
E_\nu[|g^\flat|;H_j]\le C_\delta b\), with \(C_\delta\) independent of \(b\).

This bound controls a mean of \(|g^\flat|\) under \(\nu\). Equation \eqref{eq:signed-history} also
involves the exponential under \(\mu\). The density factor \(X\) allows us to use the same bound for
both, since at arrival
\[
 X\,|e^{-g^\flat}-1|=e^{g_T+\Gamma}|1-e^{g^\flat}|.
\]
On \(H_j\) the positive part of the diagonal error has a class bound, \(|g_T|+|\Gamma|\le C/j\) and
\(|A-R_Q(n)|=O_{b,\delta}(1/n)\). We apply \(|e^y-1|\le e^{y^+}|y|\) to the density \(Xe^{A-R_Q(n)}\).
On \(F_b\setminus H_j\) the weighted probability is \(O(e^{-2j})\). For \(0<X\le1\) we have \(X|\log
X|\le1/e\), and above one we use \(X|\log X|\le X(C_0+\beta j)\) in the outer class. The contribution
of this event to the left-hand side of \eqref{eq:arrival-estimate} before scaling is
\(O((1+j)e^{-2j})\). To obtain \eqref{eq:arrival-estimate} we first fix \(\delta\), the margins,
\(\rho\) and the class constants, then \(\delta_*\) and finally the cut-off. The threshold for \(n\)
may depend on \(b\). The constant \(C_\delta\) works for all these cut-offs, as required by the final
comparison of two cut-offs in Appendix~\ref{app:closing}.
